%% parts/p07/normalization-and-curry-howard.tex --- Phase-4 P7.3
%% Migrated source: v1.1 lines 11961--12128 and 12218--12247.
%% Apparatus A-072/A-151/A-152.

\chapter{Normalization and Curry--Howard}
\label{sec:p07-normalization-and-curry-howard}
\chaptermark{Normalization and Curry--Howard}

Natural-deduction normalization and sequent cut elimination are related proof
transformations, not one theorem.  The former concerns a reduction relation
on derivation trees; the latter removes a primitive sequent rule.  The exact
natural-deduction calculus here is
\hyperref[sheet:nj]{\texttt{sheet:nj}}, and the exact typed term calculus is
\hyperref[sheet:stlc]{\texttt{sheet:stlc}}.  The normalization treatment
originates with \cite{Prawitz1965}; see also
\cite{TroelstraSchwichtenberg2000}.

\section{Reduction and natural-deduction normalization}

\begin{hbdefinition}{def:p07-abstract-reduction}{Abstract reduction relation and normalization}
An \emph{abstract reduction relation} is a binary relation
\(\longrightarrow\) on a set \(X\).  An element is \emph{normal} iff it has no
\(\longrightarrow\)-successor; it is \emph{weakly normalizing} iff some
finite reduction from it ends at a normal element, and \emph{strongly
normalizing} iff no infinite reduction sequence starts from it.  The relation
is weakly or strongly normalizing iff every element has the corresponding
property.  This terminology classifies elements and relations, not formulas
or semantic consequence.  The strong notion is nonredundant because one
terminating route does not exclude another infinite route; it is used by
Theorem~\ref{thm:p07-stlc-strong-normalization}.
\index{abstract reduction relation}
\index{normal form}
\index{weak normalization}
\index{strong normalization}
\end{hbdefinition}

\begin{hbexample}{def:p07-weak-not-strong}{Weak normalization need not be strong}
On \(X=\{a,b\}\), let \(a\longrightarrow a\) and
\(a\longrightarrow b\), with no successor of \(b\).  Both elements have a
finite route to a normal element, but the infinite constant sequence at
\(a\) witnesses failure of strong normalization.  Thus the two predicates
in Definition~\ref{def:p07-abstract-reduction} cannot be interchanged.
\end{hbexample}

\begin{hbdefinition}{def:p07-nj-normality}{NJ detours, reduction, and normal derivations}
In a \texttt{sheet:nj} derivation, the \emph{major premise} of an elimination
is the premise exhibiting the connective or quantifier eliminated.  A
\emph{detour} is a segment in which a formula is introduced and immediately
used as that major premise.  Let
\(\longrightarrow_{\mathsf{NJ}}\) be the compatible closure of the principal
reductions for matching introduction--elimination pairs and the permutation
conversions that move later eliminations above \(\lor E\) and \(\exists E\),
with bound variables alpha-renamed to preserve the sheet's discharge and
eigenvariable conditions.  A derivation is \emph{normal} iff it has no detour
and satisfies those two permutation conditions.

This relation is needed to make \emph{reduces} in the theorem exact.  It is a
metalinguistic relation on finite derivation trees and differs from term
substitution inside formulas: implication reduction substitutes a derivation
for the assumption occurrences carrying one discharge label.
\index{detour!natural deduction}
\index{major premise}
\index{normal derivation}
\end{hbdefinition}

\begin{hbmethod}{def:p07-normalization-method}{Normalization method}
\par\noindent
Reduce detours, then permute eliminations upward across disjunction and
existential eliminations.  Control the process by a
lexicographic measure whose first coordinate is the greatest rank of a
maximal formula and whose later coordinates record the lengths of maximal
segments.  A mere count of formula occurrences is inadequate because a
principal reduction can duplicate a subderivation.  The method names this
specific transformation of \(\mathsf{NJ}\) derivations; the annex enumerates
the missing calculus-specific cases.
\end{hbmethod}

\begin{hbtheorem}{thm:p07-normalization}{sketched=annex:p7.norm.1}{Normalization for intuitionistic natural deduction}
\par\noindent
Every derivation on \texttt{sheet:nj} reduces by
\(\longrightarrow_{\mathsf{NJ}}\) to a normal derivation with the same
conclusion and no undischarged assumption formula not already undischarged in
the original derivation.  Equivalently, modulo alpha-equivalence, the set of
undischarged assumption formulas in the result is a subset of the original
set.
\end{hbtheorem}

\begin{proofboundary}{def:p07-normalization-proof-boundary}{Scope of the normalization sketch}
The calculus has exactly the finite-tree rules
\(\land,\lor,\to,\bot,\forall,\exists\) on \texttt{sheet:nj}, defined
negation, total capture-avoiding formula substitution, and the displayed
nonvacuous-discharge and eigenvariable conditions.  It excludes RAA,
primitive-negation rules, equality, and induction.  The sketch omits the full
principal-reduction table, the \(\lor E/\exists E\) permutation table, and
the proof that the chosen maximal-formula/maximal-segment measure decreases.
\end{proofboundary}

\section{The simply typed lambda calculus}

\begin{hbdefinition}{def:p07-stlc}{Arrow-only simply typed lambda calculus}
Fix a set \(\mathsf{Base}\) of base types.  With type metavariables
\(\alpha,\beta\), term metavariables \(M,N\), and a finite typing context
\(\Theta_{\mathrm{ty}}\) assigning distinct term variables to types, the
types form the least set containing every \(\mathsf b\in\mathsf{Base}\) and
closed under \((\alpha,\beta)\mapsto
\alpha\mathbin{\to_{\mathrm{ty}}}\beta\).  Terms form the least set containing
every variable and closed under typed abstraction
\(\lambda x{:}\alpha.M\) and application \(MN\).
The judgment
\(\Theta_{\mathrm{ty}}\derives_{\mathsf{STLC}}M{:}\alpha\), its variable,
abstraction, and application rules, and capture-avoiding beta contraction are
exactly those on \hyperref[sheet:stlc]{\texttt{sheet:stlc}}.  The registered
type arrow \(\to_{\mathrm{ty}}\) is a metalinguistic type former, not the
object-language implication \(\to\).  This entry defines an arrow-only
calculus: products, sums, polymorphism, and dependent types are absent.
\index{simply typed lambda calculus}
\index{type!arrow}
\end{hbdefinition}

\begin{hbtheorem}{thm:p07-stlc-subject-reduction}{sketched=annex:p7.norm.2}{Subject reduction for STLC}
\par\noindent
If \(\Theta_{\mathrm{ty}}\derives_{\mathsf{STLC}}M{:}\alpha\) and
\(M\longrightarrow_\beta N\), then
\(\Theta_{\mathrm{ty}}\derives_{\mathsf{STLC}}N{:}\alpha\).
\end{hbtheorem}

\begin{hbtheorem}{thm:p07-stlc-strong-normalization}{sketched=annex:p7.norm.3}{Strong normalization for STLC}
Every term typable on \texttt{sheet:stlc} is strongly normalizing for
compatible beta reduction: no infinite beta-reduction sequence begins at a
well-typed term.
\end{hbtheorem}

\begin{hbwarning}{def:p07-normalization-distinctions}{Four properties that must remain separate}
Natural-deduction normalization concerns derivation trees; STLC subject
reduction preserves a type across one term-reduction step; strong
normalization excludes infinite term reductions; confluence compares
divergent reductions.  None of these predicates defines or implies all the
others merely because beta reduction participates in their statements.
\end{hbwarning}

\begin{hbexample}{def:p07-k-term-weakening}{The typed K term and weakening}
For types \(\alpha,\beta\), the term
\(\lambda x{:}\alpha.\lambda y{:}\beta.x\) has type
\(\alpha\mathbin{\to_{\mathrm{ty}}}
(\beta\mathbin{\to_{\mathrm{ty}}}\alpha)\).  Under Curry--Howard, the unused
variable \(y\) records antecedent weakening in the corresponding derivation
of \(\alpha\to(\beta\to\alpha)\); it is not a hidden beta reduction.
\end{hbexample}

\section{Curry--Howard and categorical semantics}

\begin{hbdefinition}{def:p07-curry-howard}{Curry--Howard setting and categorical standing}
The \emph{Curry--Howard correspondence} used here is the structural
translation between a fixed intuitionistic natural-deduction calculus and a
fixed typed term calculus: formulas map to types, assumption occurrences to
typed variables, derivations to typing derivations, discharge to abstraction,
elimination to application, and normalization to computation.  It compares
specified formal objects; it does not identify every proof with a program
without a term assignment and conversion relation.

For the third column of Table~\ref{tab:p07-curry-howard}, a
\emph{bicartesian closed category} means a category with finite products
(hence a terminal object), finite coproducts (hence an initial object), and,
for each object \(X\), a right adjoint to \(-\times X\) (hence
exponentials).  This hypothesis is necessary because products, sums, truth,
falsity, and implication have no uniform readings in an arbitrary category.
For this paragraph, \emph{beta--eta equivalence} is the least congruence on
well-typed terms containing the symmetric closure of the registered beta
contraction and the typed eta equation that identifies
\(\lambda x{:}\alpha.Mx\) with \(M\) when \(M\) has a function type with
domain \(\alpha\) and \(x\) is not free in \(M\).
It is a metalinguistic quotient used only for the categorical dictionary,
not an additional directed reduction rule on \texttt{sheet:stlc}.  The
selected syntactic category quotients terms by this equivalence; that
quotient is what makes substitution correspond to composition.  More
exactly, its objects are the selected types and a
morphism from \(\alpha\) to \(\beta\) is a beta--eta equivalence class of
judgments
\(x{:}\alpha\derives_{\mathsf{STLC}}M{:}\beta\); the variable judgment is
the identity and capture-avoiding substitution is composition.  Conversely,
the \emph{internal-language construction} used in the figure represents each
object of a chosen bicartesian closed category by a type and each morphism
from \(X\) to \(Y\) by a typed unary operation named \(f\), with the chosen
products, coproducts, and exponentials supplying the corresponding
constructors.  Both
are metalinguistic constructions relative to fixed data: the syntactic
category sends syntax to a category, whereas the internal language sends a
chosen category to syntax.  The distinction is used on the two directions of
the categorical edge in the figure.  These definitions also distinguish
term assignment from categorical semantics and are used by the table and
Figure~\ref{fig:p07-curry-howard-triangle}.
\index{Curry--Howard correspondence}
\index{bicartesian closed category}
\end{hbdefinition}

\begingroup\small
\begin{longtable}{@{}p{.24\textwidth}p{.29\textwidth}p{.38\textwidth}@{}}
\caption{Logic, type theory, and category dictionary (apparatus A-072).}
\label{tab:p07-curry-howard}\\
\toprule
Intuitionistic logic & Typed lambda calculus & Chosen bicartesian closed
category \\
\midrule
\endfirsthead
\toprule
Intuitionistic logic & Typed lambda calculus & Chosen bicartesian closed
category \\
\midrule
\endhead
Formula & Type & Object \\
Assumption context & Typed context, represented by a product of types &
Context object, represented by the corresponding categorical product \\
Derivation of \(\varphi\) from \(\Gamma\) & Term \(M\) of type \(\alpha\)
in a typed context \(\Theta_{\mathrm{ty}}\) & Morphism from the context object
to the object interpreting \(\varphi\) \\
Implication \(\varphi\to\psi\) & Function type
\(\alpha\mathbin{\to_{\mathrm{ty}}}\beta\) & Exponential \(Y^X\), where
\(X,Y\) interpret \(\varphi,\psi\) \\
Conjunction & Product type & Categorical product \\
Truth & Unit type & Terminal object \\
Disjunction & Sum type & Categorical coproduct \\
Falsity & Empty type & Initial object \\
Proof substitution / Cut & Capture-avoiding term substitution & Morphism
composition \\
Beta--eta conversion & Term equality after the declared conversions &
Equality of denoted morphisms \\
\bottomrule
\end{longtable}
\endgroup

The product, sum, unit, and empty-type rows describe the standard extension
of arrow-only \texttt{sheet:stlc}; they do not silently add constructors or
reduction rules to that sheet.  Categorical details are citation-supported by
\cite{GirardLafontTaylor1989,LambekScott1986} and receive full treatment in
the later categorical-logic part.

\begin{figure}[!htbp]
\centering
\resizebox{.9\linewidth}{!}{%
\begin{tikzpicture}[font=\small,>=Latex]
  \node[draw,rounded corners,align=center] (logic) at (0,2.8)
    {intuitionistic logic\\formulas and derivations};
  \node[draw,rounded corners,align=center] (types) at (-4.2,0)
    {typed lambda calculus\\types and terms modulo conversions};
  \node[draw,rounded corners,align=center] (cats) at (4.2,0)
    {bicartesian closed category\\objects and morphisms};
  \draw[<->] (logic) -- node[above left,align=center]
    {Curry--Howard\\term assignment} (types);
  \draw[<->] (types) -- node[below,align=center]
    {categorical semantics} (cats);
  \draw[<->] (cats) -- node[above right,align=center]
    {internal language /\\syntactic category} (logic);
\end{tikzpicture}
}
\caption{The Curry--Howard--categorical triangle (apparatus A-151).  Each edge
is a specified construction, not an equality sign; the categorical vertex
has the hypotheses of Definition~\ref{def:p07-curry-howard}.}
\label{fig:p07-curry-howard-triangle}
\end{figure}

\clearpage
\begin{landscape}
\begin{figure}[p]
\centering
\begin{minipage}[t]{.49\linewidth}
\centering\scriptsize
\textbf{STLC typing instance}\\[4pt]
\begin{tikzpicture}[font=\scriptsize,>=Latex,every node/.style={align=center}]
  \node (f) at (-2.1,4.2)
    {\(\Theta_{\mathrm{ty}}\derives_{\mathsf{STLC}}f{:}
      \beta\mathbin{\to_{\mathrm{ty}}}\gamma\)};
  \node (g) at (.2,4.2)
    {\(\Theta_{\mathrm{ty}}\derives_{\mathsf{STLC}}g{:}
      \alpha\mathbin{\to_{\mathrm{ty}}}\beta\)};
  \node (x) at (2.2,4.2)
    {\(\Theta_{\mathrm{ty}}\derives_{\mathsf{STLC}}x{:}\alpha\)};
  \node (gx) at (1.2,3.25)
    {\(\Theta_{\mathrm{ty}}\derives_{\mathsf{STLC}}gx{:}\beta\)};
  \node (fgx) at (0,2.3)
    {\(\Theta_{\mathrm{ty}}\derives_{\mathsf{STLC}}f(gx){:}\gamma\)};
  \node (lx) at (0,1.45)
    {\(f{:}\beta\mathbin{\to_{\mathrm{ty}}}\gamma,\,
      g{:}\alpha\mathbin{\to_{\mathrm{ty}}}\beta
      \derives_{\mathsf{STLC}}\lambda x{:}\alpha.f(gx){:}
      \alpha\mathbin{\to_{\mathrm{ty}}}\gamma\)};
  \node (lg) at (0,.65)
    {\(f{:}\beta\mathbin{\to_{\mathrm{ty}}}\gamma
      \derives_{\mathsf{STLC}}
      \lambda g{:}\alpha\mathbin{\to_{\mathrm{ty}}}\beta.
      \lambda x{:}\alpha.f(gx){:}
      (\alpha\mathbin{\to_{\mathrm{ty}}}\beta)
      \mathbin{\to_{\mathrm{ty}}}
      (\alpha\mathbin{\to_{\mathrm{ty}}}\gamma)\)};
  \node (lf) at (0,-.25)
    {\(\varnothing\derives_{\mathsf{STLC}}
      \lambda f{:}\beta\mathbin{\to_{\mathrm{ty}}}\gamma.
      \lambda g{:}\alpha\mathbin{\to_{\mathrm{ty}}}\beta.
      \lambda x{:}\alpha.f(gx){:}
      (\beta\mathbin{\to_{\mathrm{ty}}}\gamma)
      \mathbin{\to_{\mathrm{ty}}}
      ((\alpha\mathbin{\to_{\mathrm{ty}}}\beta)
      \mathbin{\to_{\mathrm{ty}}}
      (\alpha\mathbin{\to_{\mathrm{ty}}}\gamma))\)};
  \draw[->] (g) -- (gx);
  \draw[->] (x) -- node[right]
    {\hyperref[rule:stlc-app]{\(\to_{\mathrm{ty}}E\)}} (gx);
  \draw[->] (f) -- (fgx);
  \draw[->] (gx) -- node[right]
    {\hyperref[rule:stlc-app]{\(\to_{\mathrm{ty}}E\)}} (fgx);
  \draw[->] (fgx) -- node[right]
    {\hyperref[rule:stlc-abs]{\(\to_{\mathrm{ty}}I\)}} (lx);
  \draw[->] (lx) -- node[right]
    {\hyperref[rule:stlc-abs]{\(\to_{\mathrm{ty}}I\)}} (lg);
  \draw[->] (lg) -- node[right]
    {\hyperref[rule:stlc-abs]{\(\to_{\mathrm{ty}}I\)}} (lf);
\end{tikzpicture}
\end{minipage}\hfill
\begin{minipage}[t]{.49\linewidth}
\centering\scriptsize
\textbf{Intuitionistic ND instance}\\[4pt]
\begin{tikzpicture}[font=\scriptsize,>=Latex,every node/.style={align=center}]
  \node (f) at (-2.1,4.2) {\([\psi\to\chi]^1\)};
  \node (g) at (.2,4.2) {\([\varphi\to\psi]^2\)};
  \node (x) at (2.2,4.2) {\([\varphi]^3\)};
  \node (b) at (1.2,3.25) {\(\psi\)};
  \node (c) at (0,2.3) {\(\chi\)};
  \node (ac) at (0,1.45) {\(\varphi\to\chi\)};
  \node (gb) at (0,.65)
    {\((\varphi\to\psi)\to(\varphi\to\chi)\)};
  \node (all) at (0,-.25)
    {\((\psi\to\chi)\to
      ((\varphi\to\psi)\to(\varphi\to\chi))\)};
  \draw[->] (g) -- (b);
  \draw[->] (x) -- node[right]
    {\hyperref[rule:nj-imp-e]{\(\to E\)}} (b);
  \draw[->] (f) -- (c);
  \draw[->] (b) -- node[right]
    {\hyperref[rule:nj-imp-e]{\(\to E\)}} (c);
  \draw[->] (c) -- node[right]
    {\hyperref[rule:nj-imp-i]{\(\to I,3\)}} (ac);
  \draw[->] (ac) -- node[right]
    {\hyperref[rule:nj-imp-i]{\(\to I,2\)}} (gb);
  \draw[->] (gb) -- node[right]
    {\hyperref[rule:nj-imp-i]{\(\to I,1\)}} (all);
\end{tikzpicture}
\end{minipage}
\caption{Composition as one complete STLC typing instance beside its
intuitionistic natural-deduction instance (apparatus A-152).  Generic rules
remain on \texttt{sheet:stlc} and \texttt{sheet:nj}; arrow labels point to
their registered rows.  Here
\(\Theta_{\mathrm{ty}}=
f{:}\beta\mathbin{\to_{\mathrm{ty}}}\gamma,\,
g{:}\alpha\mathbin{\to_{\mathrm{ty}}}\beta,\,x{:}\alpha\).}
\label{fig:p07-stlc-nd-derivation}
\end{figure}
\end{landscape}
\clearpage

\paragraph{Source disposition.}
The source composition example is retained and expanded into
Figure~\ref{fig:p07-stlc-nd-derivation}.  Its typed-K exercise is retired as
an exercise but its weakening distinction survives in
Example~\ref{def:p07-k-term-weakening}.  The outermost-beta subject-reduction
exercise is retired with the exact typing-substitution and capture-avoidance
obligations listed in Annex~\ref{annex:p7.norm.2}.  The weak-versus-strong
normalization exercise is retired with both predicates and a separating
example supplied before either theorem.  No exercise solution is relabeled
as a proof.
